\chapter{Memory and Retrodiction}

The memory layer is developed downstream of Temporal Transport. The current executable branch combines the existing oriented memory coordinate with a Kepler--Newton evolution law parameterized by the internal elapsed activity variable.

\section{Oriented memory coordinate}

Assign a complex memory coordinate
\[
m(s)=x_M(s)+iy_M(s)=r_M(s)e^{i\theta_M(s)}.
\]
The oriented edge functional is
\[
\boxed{
A^{(M)}_{ab}
=\lambda_M\operatorname{Im}[m(a)^*m(b)]
=\lambda_M r_a r_b\sin(\theta_b-\theta_a).
}
\]
For a closed polygonal orbit,
\[
\sum_C A^{(M)}_e
=2\lambda_M\mathcal A_M(C),
\]
where \(\mathcal A_M(C)\) is the signed enclosed area.

\section{Internal-time Newton dynamics}

The internal evolution increment inherited from temporal activity is
\[
d\tau_{\rm int}=\frac{\mathfrak a}{\mathfrak a_\star}\,d\lambda.
\]
On the memory plane we introduce the effective central parameter \(\mu_M>0\) and
\[
V_M(r_M)=-\frac{\mu_M}{r_M}.
\]
The smooth memory trajectory obeys
\[
\boxed{
\frac{d^2m}{d\tau_{\rm int}^2}
=-\mu_M\frac{m}{|m|^3}.
}
\]
The origin and calibration of \(\mu_M\) remain an explicit derivation target.

\section{Kepler structure}

Define
\[
\varepsilon_M
=\frac12\left|\frac{dm}{d\tau_{\rm int}}\right|^2
-\frac{\mu_M}{r_M},
\qquad
h_M
=\operatorname{Im}\!\left[m^*\frac{dm}{d\tau_{\rm int}}\right].
\]
The central-force law gives the areal relation
\[
\boxed{
\frac{d\mathcal A_M}{d\tau_{\rm int}}
=\frac{h_M}{2}.
}
\]
The corresponding conic is
\[
r_M(\nu)
=\frac{p_M}{1+e_M\cos\nu},
\qquad
p_M=\frac{h_M^2}{\mu_M}.
\]
For a bound nondegenerate orbit,
\[
a_M=-\frac{\mu_M}{2\varepsilon_M},
\qquad
T_M^2=\frac{4\pi^2}{\mu_M}a_M^3.
\]
Combining the areal law with the earlier memory circulation yields
\[
\boxed{
\frac{d}{d\tau_{\rm int}}
\left(\sum_C A_e^{(M)}\right)
=\lambda_M h_M.
}
\]

\section{Bifurcation impulse}

A localized event can modify the orbital state through
\[
\boxed{
\mathbf v_M^+
=\mathbf v_M^-+\Delta\mathbf v_M.
}
\]
The event-imprint projection that generates \(\Delta\mathbf v_M\) is kept as a separate derivation target. The resulting orbital state can be classified by \(\varepsilon_M\), \(h_M\), eccentricity and the conic branch.

The current numerical reference uses velocity--Verlet propagation and records targeted structural tests in \texttt{validation/KEPLER\_MEMORY\_DYNAMICS\_V0\_1.json}. Retrodiction opens after the memory update law receives its own admission receipt.
